\chapter{Hamiltonian Time Evolution, Interference, and the Schr\"odinger Equation}

\section{The Classical Hamiltonian Generates a Flow}

We now begin the serious study of change with time. It is useful to start with
the classical picture, not because we shall need every detail of classical
mechanics, but because one structural idea survives intact: the Hamiltonian is
the generator of time evolution.

For a system with finitely many classical states, a law of evolution may be
drawn as arrows from one state to the next. A coin could pass through a
sequence such as heads, tails, heads, and so on. Since there is no continuous
path between two isolated states, this picture naturally describes updates at
discrete sampling times.

Continuous classical motion requires a continuous state space. A particle on
a line has, at each instant, a position \(x\) and momentum \(p\). Its state is
therefore a point in the two-dimensional phase space with coordinates
\((x,p)\). As time passes, that point follows a flow line.

The flow is determined by a function of phase space called the Hamiltonian.
For a particle of mass \(m\) in a potential \(V(x)\), the standard example is
\begin{equation}
  H(x,p)=\frac{p^2}{2m}+V(x).
  \label{eq:l8-classical-H}
\end{equation}
The momentum-dependent term is the kinetic energy and the
position-dependent term is the potential energy. For many particles there is
one momentum conjugate to each coordinate, but one pair is enough to display
the structure.

\begin{figure}[tbp]
  \centering
  \includegraphics[width=0.70\textwidth]{lecture_08_figure_02.png}
  \caption{The classical Hamiltonian for a particle on a line, with the
  kinetic and potential terms separated.}
  \lectureframe{00:05:22}
\end{figure}

Hamilton's equations turn this energy function into motion:
\begin{equation}
  \dot x=\frac{\partial H}{\partial p},
  \qquad
  \dot p=-\frac{\partial H}{\partial x}.
  \label{eq:l8-Hamilton-equations}
\end{equation}
For Equation~\eqref{eq:l8-classical-H},
\(\partial H/\partial p=p/m\), which is the velocity. The second equation
states that the rate of change of momentum is the force,
\(-\partial V/\partial x\).

Now let \(F(x,p)\) be any classical observable. The chain rule gives
\begin{align}
  \frac{dF}{dt}
  &=
  \frac{\partial F}{\partial x}\dot x
  +\frac{\partial F}{\partial p}\dot p\\
  &=
  \frac{\partial F}{\partial x}
  \frac{\partial H}{\partial p}
  -
  \frac{\partial F}{\partial p}
  \frac{\partial H}{\partial x}.
  \label{eq:l8-classical-chain}
\end{align}
The combination on the right is the Poisson bracket:
\begin{equation}
  \{F,H\}_{\mathrm{PB}}
  =
  \frac{\partial F}{\partial x}
  \frac{\partial H}{\partial p}
  -
  \frac{\partial F}{\partial p}
  \frac{\partial H}{\partial x}.
  \label{eq:l8-Poisson}
\end{equation}
Thus
\begin{equation}
  \boxed{\frac{dF}{dt}=\{F,H\}_{\mathrm{PB}}.}
  \label{eq:l8-classical-evolution}
\end{equation}

\begin{figure}[tbp]
  \centering
  \includegraphics[width=0.76\textwidth]{lecture_08_figure_03.png}
  \caption{Hamilton's equations inserted into the chain rule produce the
  Poisson-bracket law for an arbitrary observable.}
  \lectureframe{00:08:38}
\end{figure}

One need not memorize the bracket at this stage. Its message is more
important than its detailed form: once \(H\) is known, every phase-space
quantity has a determined time derivative. The Hamiltonian generates the
flow.

\section{Quantum Evolution Preserves Relations Between States}

The quantum statement begins from a different picture. A state is a vector,
and the physically meaningful relation between two normalized states is their
inner product. Orthogonal states can be distinguished with certainty;
nonorthogonal states have a nonzero transition probability. Closed-system
evolution must preserve all of these relations.

Suppose that after an elapsed time \(T\),
\begin{equation}
  \ket{A}\longmapsto\hat U(T)\ket{A},
  \qquad
  \ket{B}\longmapsto\hat U(T)\ket{B}.
  \label{eq:l8-evolution-map}
\end{equation}
The corresponding transformed bra is
\begin{equation}
  \bra{B}\longmapsto\bra{B}\hat U^\dagger(T).
\end{equation}
Preservation of the inner product requires
\begin{align}
  \braket{B}{A}
  &=
  \bra{B}\hat U^\dagger(T)\hat U(T)\ket{A}.
\end{align}
This must hold for every pair \(\ket A,\ket B\). Therefore
\begin{equation}
  \boxed{\hat U^\dagger(T)\hat U(T)=\mathbf 1.}
  \label{eq:l8-unitary-condition}
\end{equation}
In finite dimensions this also implies
\(\hat U(T)\hat U^\dagger(T)=\mathbf1\). The evolution operator is unitary.

\begin{figure}[tbp]
  \centering
  \includegraphics[width=0.72\textwidth]{lecture_08_figure_04.png}
  \caption{The common evolution of two states preserves their inner product
  and leads to \(\hat U^\dagger\hat U=\mathbf1\).}
  \lectureframe{00:14:20}
\end{figure}

There is a useful geometric image here. The vectors move, but the whole state
space moves rigidly: norms and mutual angles remain fixed. In particular,
orthogonal states remain orthogonal. For the linear evolution considered
here, this preservation law is the governing principle.

Writing the state at an arbitrary initial time \(t\), the evolution law is
\begin{equation}
  \ket{\psi(t+T)}
  =
  \hat U(T)\ket{\psi(t)}.
  \label{eq:l8-finite-evolution}
\end{equation}
No elapsed time means no change,
\begin{equation}
  \hat U(0)=\mathbf1.
  \label{eq:l8-U-zero}
\end{equation}
Successive intervals compose:
\begin{equation}
  \hat U(T_2)\hat U(T_1)=\hat U(T_1+T_2)
  \label{eq:l8-group-law}
\end{equation}
when the Hamiltonian has no explicit time dependence. This is why a long
evolution can be assembled from many short steps.

\section{The Infinitesimal Generator}

Let \(\epsilon\) be short enough that terms of order \(\epsilon^2\) may be
neglected. Since \(\hat U(0)=\mathbf1\), the most general first-order form can
be written
\begin{equation}
  \hat U(\epsilon)
  =
  \mathbf1-\frac{i\epsilon}{\hbar}\hat H
  +O(\epsilon^2).
  \label{eq:l8-infinitesimal-U}
\end{equation}
At first, \(\hat H\) is simply the name of an operator. The factors
\(-i/\hbar\) are chosen so that \(\hat H\) has units of energy and so that the
classical correspondence will later take its standard form.

\begin{figure}[tbp]
  \centering
  \includegraphics[width=0.74\textwidth]{lecture_08_figure_05.png}
  \caption{A short-time evolution differs from the identity by a term
  proportional to \(-i\epsilon\hat H/\hbar\).}
  \lectureframe{00:19:54}
\end{figure}

Taking the adjoint gives
\begin{equation}
  \hat U^\dagger(\epsilon)
  =
  \mathbf1+\frac{i\epsilon}{\hbar}\hat H^\dagger
  +O(\epsilon^2).
\end{equation}
Unitarity through first order then says
\begin{align}
  \mathbf1
  &=
  \hat U^\dagger(\epsilon)\hat U(\epsilon)\\
  &=
  \left(
    \mathbf1+\frac{i\epsilon}{\hbar}\hat H^\dagger
  \right)
  \left(
    \mathbf1-\frac{i\epsilon}{\hbar}\hat H
  \right)
  +O(\epsilon^2)\\
  &=
  \mathbf1
  +\frac{i\epsilon}{\hbar}
  \left(\hat H^\dagger-\hat H\right)
  +O(\epsilon^2).
\end{align}
The linear term must vanish, so
\begin{equation}
  \boxed{\hat H^\dagger=\hat H.}
  \label{eq:l8-H-Hermitian}
\end{equation}

\begin{figure}[tbp]
  \centering
  \includegraphics[width=0.74\textwidth]{lecture_08_figure_06.png}
  \caption{Expanding \(\hat U^\dagger\hat U=\mathbf1\) to first order leaves
  the condition \(\hat H^\dagger-\hat H=0\).}
  \lectureframe{00:22:48}
\end{figure}

The generator of continuous unitary evolution is Hermitian. Its eigenvalues
are real and, in finite dimensions, its eigenvectors may be chosen as an
orthonormal basis. We identify this observable as the Hamiltonian and call its
eigenvalues energies. At this point that is an identification within quantum
mechanics; recovering familiar classical energy requires a further
correspondence argument.

For a time-independent Hamiltonian, compounding the short-time step gives
\begin{equation}
  \hat U(T)
  =
  \exp\left(-\frac{i}{\hbar}\hat H T\right).
  \label{eq:l8-exponential-U}
\end{equation}
The exponential is unitary precisely because \(\hat H\) is Hermitian.

\begin{classroomqa}
\audiencequestion{If \(\hbar\) appears in the denominator, what happens in
the classical limit \(\hbar\to0\)?
\lecturetimestamp{00:26:41--00:28:07}}
\lecturerresponse{The classical limit is not obtained by setting a denominator
to zero inside one fixed microscopic formula. What matters are
dimensionless actions measured in units of \(\hbar\), such as
\(\epsilon E/\hbar\). Classical behavior emerges when the relevant actions
are large compared with \(\hbar\) and quantum phases combine by stationary
phase or decoherence. The operator \(\hat H\) itself has a well-defined
classical counterpart even though the phases become rapidly varying.}
\end{classroomqa}

\section{The Abstract Schr\"odinger Equation}

Apply Equation~\eqref{eq:l8-infinitesimal-U} to a state:
\begin{equation}
  \ket{\psi(t+\epsilon)}
  =
  \left(
    \mathbf1-\frac{i\epsilon}{\hbar}\hat H
  \right)
  \ket{\psi(t)}
  +O(\epsilon^2).
  \label{eq:l8-small-state-step}
\end{equation}
Subtract \(\ket{\psi(t)}\), divide by \(\epsilon\), and take
\(\epsilon\to0\):
\begin{equation}
  \frac{\partial}{\partial t}\ket{\psi(t)}
  =
  -\frac{i}{\hbar}\hat H\ket{\psi(t)}.
\end{equation}
Multiplication by \(i\hbar\) gives
\begin{equation}
  \boxed{
  i\hbar\frac{\partial}{\partial t}\ket{\psi(t)}
  =
  \hat H\ket{\psi(t)}.}
  \label{eq:l8-abstract-Schrodinger}
\end{equation}

\begin{figure}[tbp]
  \centering
  \includegraphics[width=0.76\textwidth]{lecture_08_figure_07.png}
  \caption{The finite difference from \(t\) to \(t+\epsilon\) becomes the
  abstract Schr\"odinger equation.}
  \lectureframe{00:32:03}
\end{figure}

This is the general law of quantum time evolution. It is abstract because
\(\ket{\psi}\) may describe polarization, spin, a particle in space, or any
other quantum system. Schr\"odinger's original wave equation will appear only
after we choose a coordinate representation and specify \(\hat H\).

There is a useful parallel with momentum. In the coordinate representation,
\begin{equation}
  \hat p=-i\hbar\frac{\partial}{\partial x}
\end{equation}
relates spatial translation to momentum, while
Equation~\eqref{eq:l8-abstract-Schrodinger} relates time translation to
energy. The sign conventions are conventional; their consistent use is what
matters.

\section{A Polarization Doublet in an Anisotropic Medium}

\subsection{The Energy Eigenstates}

Consider a photon of fixed wavelength moving through a material whose two
preferred polarization axes have different energies. Choose those material
axes as \(x\) and \(y\), and suppose
\begin{equation}
  \hat H\ket{x}=E_1\ket{x},
  \qquad
  \hat H\ket{y}=E_2\ket{y}.
  \label{eq:l8-polarization-eigenstates}
\end{equation}
Here \(x\) and \(y\) label polarization, not position. The assumption is that
the material itself distinguishes these two directions.

\begin{figure}[tbp]
  \centering
  \includegraphics[width=0.70\textwidth]{lecture_08_figure_08.png}
  \caption{The material assigns energies \(E_1\) and \(E_2\) to its two
  orthogonal polarization eigenstates.}
  \lectureframe{00:35:54}
\end{figure}

A general polarization state is
\begin{equation}
  \ket{\psi(t)}
  =
  \alpha(t)\ket{x}+\beta(t)\ket{y}
  =
  \begin{pmatrix}
    \alpha(t)\\
    \beta(t)
  \end{pmatrix}.
  \label{eq:l8-polarization-state}
\end{equation}
The probabilities of obtaining \(x\) and \(y\) polarization are
\(|\alpha(t)|^2\) and \(|\beta(t)|^2\).

In the \((\ket{x},\ket{y})\) basis,
\begin{equation}
  \hat H
  =
  \begin{pmatrix}
    E_1&0\\
    0&E_2
  \end{pmatrix}.
  \label{eq:l8-diagonal-H}
\end{equation}
The diagonal order follows directly from
Equation~\eqref{eq:l8-polarization-eigenstates}.

\begin{figure}[tbp]
  \centering
  \includegraphics[width=0.72\textwidth]{lecture_08_figure_09.png}
  \caption{The polarization eigenvalue equations used to construct the
  diagonal Hamiltonian.}
  \lectureframe{00:39:31}
\end{figure}

The Schr\"odinger equation becomes
\begin{equation}
  i\hbar
  \frac{d}{dt}
  \begin{pmatrix}
    \alpha\\
    \beta
  \end{pmatrix}
  =
  \begin{pmatrix}
    E_1&0\\
    0&E_2
  \end{pmatrix}
  \begin{pmatrix}
    \alpha\\
    \beta
  \end{pmatrix}.
\end{equation}
The two component equations decouple:
\begin{equation}
  i\hbar\dot\alpha=E_1\alpha,
  \qquad
  i\hbar\dot\beta=E_2\beta.
  \label{eq:l8-component-equations}
\end{equation}
Their solutions are
\begin{equation}
  \alpha(t)=\alpha(0)e^{-iE_1t/\hbar},
  \qquad
  \beta(t)=\beta(0)e^{-iE_2t/\hbar}.
  \label{eq:l8-component-solutions}
\end{equation}
Thus
\begin{equation}
  \ket{\psi(t)}
  =
  \begin{pmatrix}
    \alpha(0)e^{-iE_1t/\hbar}\\
    \beta(0)e^{-iE_2t/\hbar}
  \end{pmatrix}.
  \label{eq:l8-evolved-polarization}
\end{equation}

\begin{figure}[tbp]
  \centering
  \includegraphics[width=0.70\textwidth]{lecture_08_figure_10.png}
  \caption{Each energy component accumulates its own phase,
  \(e^{-iE_jt/\hbar}\).}
  \lectureframe{00:43:55}
\end{figure}

If we measure in the energy basis itself, nothing seems to change:
\begin{equation}
  |\alpha(t)|^2=|\alpha(0)|^2,
  \qquad
  |\beta(t)|^2=|\beta(0)|^2.
  \label{eq:l8-fixed-energy-probabilities}
\end{equation}
Each exponential has unit modulus. The relative phase is invisible to an
\(x/y\) analyzer, but it is not physically irrelevant.

\subsection{Interference in a Rotated Analyzer}

Let
\begin{equation}
  \ket{d_+}
  =
  \frac{1}{\sqrt2}
  \begin{pmatrix}1\\1\end{pmatrix},
  \qquad
  \ket{d_-}
  =
  \frac{1}{\sqrt2}
  \begin{pmatrix}1\\-1\end{pmatrix}
  \label{eq:l8-diagonal-polarizations}
\end{equation}
be the \(+45^\circ\) and \(-45^\circ\) linear polarizations. The amplitude
to pass a \(+45^\circ\) analyzer after time \(t\) is
\begin{equation}
  \braket{d_+}{\psi(t)}
  =
  \frac{1}{\sqrt2}
  \left[
    \alpha(0)e^{-iE_1t/\hbar}
    +\beta(0)e^{-iE_2t/\hbar}
  \right].
  \label{eq:l8-diagonal-amplitude-general}
\end{equation}

\begin{figure}[tbp]
  \centering
  \includegraphics[width=0.74\textwidth]{lecture_08_figure_11.png}
  \caption{Projection of the evolved two-component state onto the
  \(45^\circ\) analyzer.}
  \lectureframe{00:49:24}
\end{figure}

Prepare the photon with a first \(+45^\circ\) polarizer. Then
\begin{equation}
  \alpha(0)=\beta(0)=\frac{1}{\sqrt2},
\end{equation}
and
\begin{equation}
  \mathcal A_+(t)
  =
  \frac12
  \left(
    e^{-iE_1t/\hbar}
    +e^{-iE_2t/\hbar}
  \right).
  \label{eq:l8-diagonal-amplitude}
\end{equation}
The transmission probability is the absolute square:
\begin{align}
  P_+(t)
  &=
  |\mathcal A_+(t)|^2\\
  &=
  \frac14
  \left(
    e^{-iE_1t/\hbar}
    +e^{-iE_2t/\hbar}
  \right)
  \left(
    e^{iE_1t/\hbar}
    +e^{iE_2t/\hbar}
  \right)\\
  &=
  \frac14
  \left[
    2
    +e^{i(E_2-E_1)t/\hbar}
    +e^{-i(E_2-E_1)t/\hbar}
  \right].
  \label{eq:l8-interference-expansion}
\end{align}

\begin{figure}[tbp]
  \centering
  \includegraphics[width=0.76\textwidth]{lecture_08_figure_12.png}
  \caption{The absolute square contains two direct terms and two
  phase-sensitive cross terms.}
  \lectureframe{00:54:54}
\end{figure}

Define the energy splitting
\begin{equation}
  \Delta E=E_2-E_1.
\end{equation}
Using \(e^{i\delta}+e^{-i\delta}=2\cos\delta\), we obtain
\begin{align}
  P_+(t)
  &=
  \frac12
  \left[
    1+\cos\left(\frac{\Delta E\,t}{\hbar}\right)
  \right]\\
  &=
  \cos^2\left(\frac{\Delta E\,t}{2\hbar}\right).
  \label{eq:l8-diagonal-probability}
\end{align}
The complementary analyzer has
\begin{equation}
  P_-(t)
  =
  \sin^2\left(\frac{\Delta E\,t}{2\hbar}\right),
  \qquad
  P_+(t)+P_-(t)=1.
  \label{eq:l8-complementary-probability}
\end{equation}

\begin{figure}[tbp]
  \centering
  \includegraphics[width=0.70\textwidth]{lecture_08_figure_13.png}
  \caption{The probability for the second parallel analyzer oscillates
  between one and zero as the relative phase advances.}
  \lectureframe{00:58:44}
\end{figure}

At \(t=0\), two adjacent parallel analyzers transmit with certainty. When
\(\Delta E\,t/\hbar=\pi\), the state is orthogonal to its initial
\(+45^\circ\) polarization and passes the \(-45^\circ\) analyzer with
certainty.

Between those two times the state is generally not linearly polarized.
Removing a common phase from
Equation~\eqref{eq:l8-evolved-polarization} gives
\begin{equation}
  \ket{\psi(t)}
  \sim
  \frac{1}{\sqrt2}
  \begin{pmatrix}
    1\\
    e^{-i\Delta E t/\hbar}
  \end{pmatrix}.
  \label{eq:l8-relative-polarization-phase}
\end{equation}
It begins at \(+45^\circ\) linear polarization, becomes circular when the
relative phase is \(\pm\pi/2\), reaches \(-45^\circ\) linear polarization at
\(\pi\), and passes through elliptical states in between. The analyzer
oscillation is therefore an interference measurement of a continuously
changing polarization state.

\subsection{Only Energy Differences Are Observable}

If a constant \(C\) is added to every energy,
\begin{equation}
  \hat H\longmapsto\hat H+C\mathbf1,
\end{equation}
then
\begin{equation}
  \hat U(t)
  \longmapsto
  e^{-iCt/\hbar}\hat U(t).
\end{equation}
Every state acquires the same overall phase, which changes no transition
probability. This is why the interference depends on \(\Delta E\), not on the
arbitrary zero of energy.

\begin{classroomqa}
\audiencequestion{Must the material fill the entire region between the two
polarizers?\lecturetimestamp{01:00:52--01:01:27}}
\lecturerresponse{The idealized calculation assumes that the photon spends a
known time in a uniform anisotropic region. A finite slab works just as well:
replace \(t\) by the transit time through that slab. Outside it, any common
vacuum phase does not alter the polarization probabilities.}
\end{classroomqa}

\begin{classroomqa}
\audiencequestion{The graph uses time on its horizontal axis. How does that
become a downstream position in the experiment?
\lecturetimestamp{01:01:29--01:02:43}}
\lecturerresponse{If the photon travels through the material with group
velocity \(v\), a slab length \(L\) corresponds to an interaction time
\(t=L/v\). Moving the second analyzer downstream, or changing the material
thickness, therefore changes the accumulated relative phase
\(\Delta E L/(\hbar v)\).}
\end{classroomqa}

\begin{classroomqa}
\audiencequestion{What happens if the incident photon is circularly polarized
instead of linearly polarized?\lecturetimestamp{01:02:44--01:03:40}}
\lecturerresponse{A circular state is an equal-amplitude \(x/y\) superposition
with relative phase \(\pm\pi/2\). The material adds
\(-\Delta E t/\hbar\) to that phase, so the state moves through elliptical
polarizations, becomes linearly polarized at appropriate times, and later
reaches the opposite circular handedness. The same interference calculation
applies after changing the initial coefficients.}
\end{classroomqa}

\begin{classroomqa}
\audiencequestion{If the labels \(x\) and \(y\) are arbitrary, why should a
different choice change the prediction?
\lecturetimestamp{01:03:43--01:05:15}}
\lecturerresponse{The labels are conventional, but the axes are not. The
material has two physical principal axes. We chose coordinates aligned with
them so that \(\hat H\) is diagonal. Rotating only the coordinate labels
changes no prediction; rotating the polarizers relative to the material
changes the preparation and measurement basis and therefore reveals the
relative phase.}
\end{classroomqa}

\section{Expectation Values and Commutators}

Let \(\hat K\) be a fixed observable with no explicit time dependence. Its
expectation value in the evolving state is
\begin{equation}
  \expect{\hat K}
  =
  \bra{\psi(t)}\hat K\ket{\psi(t)}.
  \label{eq:l8-K-expectation}
\end{equation}
Only the bra and ket vary. The product rule gives
\begin{equation}
  \frac{d}{dt}\expect{\hat K}
  =
  \bra{\dot\psi}\hat K\ket{\psi}
  +
  \bra{\psi}\hat K\ket{\dot\psi}.
  \label{eq:l8-expectation-product-rule}
\end{equation}

The ket equation and its adjoint are
\begin{equation}
  \ket{\dot\psi}
  =
  -\frac{i}{\hbar}\hat H\ket{\psi},
  \qquad
  \bra{\dot\psi}
  =
  \frac{i}{\hbar}\bra{\psi}\hat H,
  \label{eq:l8-bra-ket-Schrodinger}
\end{equation}
where Hermiticity of \(\hat H\) has been used.

\begin{figure}[tbp]
  \centering
  \includegraphics[width=0.68\textwidth]{lecture_08_figure_14.png}
  \caption{The ket Schr\"odinger equation and the corresponding bra equation
  carry opposite signs.}
  \lectureframe{01:11:45}
\end{figure}

Substitution gives
\begin{align}
  \frac{d}{dt}\expect{\hat K}
  &=
  \frac{i}{\hbar}
  \bra{\psi}\hat H\hat K\ket{\psi}
  -
  \frac{i}{\hbar}
  \bra{\psi}\hat K\hat H\ket{\psi}\\
  &=
  \frac{i}{\hbar}
  \expect{[\hat H,\hat K]}\\
  &=
  -\frac{i}{\hbar}
  \expect{[\hat K,\hat H]}.
  \label{eq:l8-Heisenberg-expectation}
\end{align}

\begin{figure}[tbp]
  \centering
  \includegraphics[width=0.76\textwidth]{lecture_08_figure_15.png}
  \caption{Differentiating the expectation-value sandwich produces the two
  operator orderings \(\hat H\hat K\) and \(\hat K\hat H\).}
  \lectureframe{01:13:20}
\end{figure}

\begin{classroomqa}
\audiencequestion{Is the result written with
\([\hat K,\hat H]\) or \([\hat H,\hat K]\), and which sign belongs in front?
\lecturetimestamp{01:14:46--01:15:18}}
\lecturerresponse{Both forms are equivalent because
\([\hat H,\hat K]=-[\hat K,\hat H]\). With
\([\hat K,\hat H]=\hat K\hat H-\hat H\hat K\),
\[
  \frac{d}{dt}\expect{\hat K}
  =
  -\frac{i}{\hbar}\expect{[\hat K,\hat H]}.
\]
Writing the commutator in the opposite order changes the prefactor to
\(+i/\hbar\).}
\end{classroomqa}

\begin{figure}[tbp]
  \centering
  \includegraphics[width=0.76\textwidth]{lecture_08_figure_16.png}
  \caption{The completed commutator form of the expectation-value evolution
  law.}
  \lectureframe{01:16:37}
\end{figure}

This is the quantum counterpart of
Equation~\eqref{eq:l8-classical-evolution}. In classical mechanics, take a
Poisson bracket with \(H\); in quantum mechanics, take a commutator with
\(\hat H\). The precise correspondence is
\begin{equation}
  \{F,H\}_{\mathrm{PB}}
  \quad\longleftrightarrow\quad
  \frac{1}{i\hbar}[\hat F,\hat H]
\end{equation}
in the appropriate classical limit.

\section{Conserved Quantities}

If
\begin{equation}
  [\hat K,\hat H]=0,
  \label{eq:l8-commuting-conserved}
\end{equation}
then Equation~\eqref{eq:l8-Heisenberg-expectation} gives
\begin{equation}
  \frac{d}{dt}\expect{\hat K}=0
\end{equation}
in every state. More strongly, the spectral projectors of \(\hat K\) commute
with \(\hat H\), so the full probability distribution of \(K\) is
time-independent.

\begin{figure}[tbp]
  \centering
  \includegraphics[width=0.72\textwidth]{lecture_08_figure_17.png}
  \caption{A vanishing commutator with the Hamiltonian is the criterion for a
  time-independent observable distribution.}
  \lectureframe{01:18:49}
\end{figure}

The converse requires the same scope: if the expectation value is constant
for every state, and \(\hat K\) has no explicit time dependence, then
\([\hat K,\hat H]=0\). Constancy in one specially chosen state alone would not
be enough.

The Hamiltonian commutes with itself:
\begin{equation}
  [\hat H,\hat H]=0.
\end{equation}
Therefore a time-independent Hamiltonian is conserved. Since it represents
energy, this is energy conservation. The corresponding classical statement
is
\begin{equation}
  \{H,H\}_{\mathrm{PB}}=0.
\end{equation}
If the Hamiltonian depends explicitly on time, an additional term
\(\expect{\partial\hat H/\partial t}\) appears and energy need not be
conserved. The lecture's argument assumes no such explicit dependence.

\section{A Free Particle on a Line}

\subsection{From an Abstract Ket to a Wave Function}

For a particle on a line, the state may be represented by
\(\psi(x,t)\). Its position probability density is
\begin{equation}
  \rho(x,t)=|\psi(x,t)|^2.
\end{equation}
The Fourier transform \(\widetilde\psi(p,t)\) gives the momentum
representation, with \(|\widetilde\psi(p,t)|^2\) as the momentum probability
density. These are not new kinds of state; they are coordinate
representations of the same abstract vector.

Multiplication by \(x\) and differentiation with respect to \(x\) are linear
operators on the vector space of wave functions:
\begin{equation}
  \hat x\psi=x\psi,
  \qquad
  \hat p\psi
  =
  -i\hbar\frac{\partial\psi}{\partial x}.
  \label{eq:l8-x-p-operators}
\end{equation}

\begin{figure}[tbp]
  \centering
  \includegraphics[width=0.72\textwidth]{lecture_08_figure_18.png}
  \caption{The wave-function Schr\"odinger equation is paired with the
  differential representation of momentum.}
  \lectureframe{01:24:03}
\end{figure}

For a force-free particle, classical mechanics suggests
\begin{equation}
  H=\frac{p^2}{2m}.
\end{equation}
This does not prove the quantum rule, but it gives the correspondence
postulate
\begin{equation}
  \hat H
  =
  \frac{\hat p^2}{2m}.
  \label{eq:l8-free-H}
\end{equation}
Acting twice with momentum,
\begin{align}
  \hat p^2\psi
  &=
  \left(
    -i\hbar\frac{\partial}{\partial x}
  \right)
  \left(
    -i\hbar\frac{\partial\psi}{\partial x}
  \right)\\
  &=
  -\hbar^2
  \frac{\partial^2\psi}{\partial x^2}.
\end{align}
Consequently
\begin{equation}
  \hat H
  =
  -\frac{\hbar^2}{2m}
  \frac{\partial^2}{\partial x^2}.
  \label{eq:l8-free-H-coordinate}
\end{equation}

\begin{classroomqa}
\audiencequestion{Where do the minus sign and the powers of \(\hbar\) go in
the free-particle equation?
\lecturetimestamp{01:27:15--01:29:09}}
\lecturerresponse{Keep the abstract equation and momentum operator visible at
the same time:
\[
  i\hbar\,\partial_t\psi=\hat H\psi,
  \qquad
  \hat H=\frac{\hat p^2}{2m},
  \qquad
  \hat p=-i\hbar\,\partial_x.
\]
Squaring \(\hat p\) produces
\((-i)^2\hbar^2\partial_x^2=-\hbar^2\partial_x^2\). No
additional factor of \(i\) belongs to \(\hat H\).}
\end{classroomqa}

Substitution into the abstract equation yields Schr\"odinger's free-particle
wave equation:
\begin{equation}
  \boxed{
  i\hbar\frac{\partial\psi(x,t)}{\partial t}
  =
  -\frac{\hbar^2}{2m}
  \frac{\partial^2\psi(x,t)}{\partial x^2}.}
  \label{eq:l8-free-Schrodinger}
\end{equation}

\begin{figure}[tbp]
  \centering
  \includegraphics[width=0.74\textwidth]{lecture_08_figure_19.png}
  \caption{The complete free-particle Schr\"odinger equation with its
  \(\hbar\), \(i\), and minus-sign factors restored.}
  \lectureframe{01:30:44}
\end{figure}

The Hamiltonian is Hermitian on an appropriate domain because \(\hat p\) is
Hermitian there and the square of a Hermitian operator is Hermitian. Boundary
conditions and domains matter in a full treatment; on the line, wave
functions are taken to decay sufficiently or are handled with the standard
self-adjoint momentum domain.

\subsection{A Momentum Eigenstate}

Look for a solution that remains a momentum eigenstate:
\begin{equation}
  \psi_p(x,t)
  =
  f(t)e^{ipx/\hbar}.
  \label{eq:l8-plane-wave-ansatz}
\end{equation}
The number \(p\) is the momentum eigenvalue, because
\begin{equation}
  -i\hbar\frac{\partial}{\partial x}
  e^{ipx/\hbar}
  =
  p\,e^{ipx/\hbar}.
\end{equation}
The function \(f(t)\) is a spectator under spatial differentiation and
contains all time dependence.

\begin{figure}[tbp]
  \centering
  \includegraphics[width=0.46\textwidth]{lecture_08_figure_20.png}
  \caption{The spatial factor \(e^{ipx/\hbar}\) is an eigenfunction of the
  momentum operator.}
  \lectureframe{01:32:24}
\end{figure}

Substituting Equation~\eqref{eq:l8-plane-wave-ansatz} into
Equation~\eqref{eq:l8-free-Schrodinger} gives
\begin{equation}
  i\hbar\dot f(t)
  =
  \frac{p^2}{2m}f(t).
\end{equation}
Thus
\begin{equation}
  f(t)
  =
  f(0)
  \exp\left(
    -\frac{i}{\hbar}\frac{p^2}{2m}t
  \right),
\end{equation}
and the full solution is
\begin{equation}
  \boxed{
  \psi_p(x,t)
  =
  C
  \exp\left[
    \frac{i}{\hbar}
    \left(
      px-\frac{p^2}{2m}t
    \right)
  \right].}
  \label{eq:l8-free-plane-wave}
\end{equation}

\begin{figure}[tbp]
  \centering
  \includegraphics[width=0.46\textwidth]{lecture_08_figure_21.png}
  \caption{The time factor and spatial momentum eigenfunction combine into
  the free plane wave.}
  \lectureframe{01:36:59}
\end{figure}

\enlargethispage{\baselineskip}
The energy is \(E=p^2/(2m)\), so the familiar phase is
\((px-Et)/\hbar\). A single ideal momentum eigenstate changes only by an
energy phase and retains definite momentum. Such a plane wave is
delta-normalized rather than square-integrable; a physical localized state is
a superposition of these waves. Its different momentum components acquire
different phases, producing motion and spreading, which is the next step in
the analysis.

Finally,
\begin{equation}
  [\hat p,\hat H]
  =
  \left[
    \hat p,\frac{\hat p^2}{2m}
  \right]
  =0.
\end{equation}
Momentum is therefore conserved for a free particle. The classical
Hamiltonian flow, unitary state evolution, polarization interference,
commutator conservation law, and free-particle wave equation are all
different expressions of the same organizing principle: the Hamiltonian
generates change with time.
